\documentclass[a4paper, 10pt]{article}
%\usepackage{fontspec}
%\setmainfont{Lato}
\usepackage{pgf}
\usepackage{eurosym}
\usepackage{graphicx}
\usepackage{wasysym}
\usepackage{hyperref}
\usepackage{listings}
\usepackage{pxfonts}
\usepackage{verbatim}
\usepackage{color}
\usepackage{xcolor}
\usepackage{wrapfig}
\usepackage{enumitem}
\usepackage{booktabs}
\usepackage{gensymb}
\usepackage{tabularx}
\usepackage{currfile}

\hypersetup{
    bookmarks=true,         % show bookmarks bar?
    unicode=true,          % non-Latin characters in Acrobat’s bookmarks
    pdftoolbar=true,        % show Acrobat’s toolbar?
    pdfmenubar=true,        % show Acrobat’s menu?
    pdffitwindow=true,     % window fit to page when opened
    pdftitle={Assessments},    % title
    pdfauthor={Paul Vesey},     % author
    pdfsubject={Advanced Graphics Assignment },   % subject of the document
    pdfcreator={},   % creator of the document
    pdfproducer={xelatex}, % producer of the document
    pdfkeywords={'Graphics' }, % list of keywords
    pdfnewwindow=true,      % links in new PDF window
    colorlinks=true,       % false: boxed links; true: colored links
    linkcolor=violet,          % color of internal links (change box color with linkbordercolor)
    citecolor=magenta,        % color of links to bibliography
    filecolor=red,      % color of file links
    urlcolor=blue           % color of external links
}

\setlength\parindent{0pt}
\begin{document}

\lstset{language=HTML,
				basicstyle=\small,
				breaklines=true,
        numbers=left,
        numberstyle=\tiny,
        showstringspaces=false,
        aboveskip=-20pt,
        frame=leftline
        }
				
\begin{table}%
	\begin{minipage}{0.4\textwidth}%
			\includegraphics[width=1\textwidth]{./img/LITlogo.jpg}
	\end{minipage}
	\qquad
	\centering
	\parbox{0.4\textwidth}{
		\begin{large}			
			\begin{tabular}{| r | l |} \hline
				Subject: & \textbf{Advanced Graphics}\\
								 & \textbf{\& Visualisation}\\
				Course: & \textbf{Interior Design Y3}\\
				Session: & \textbf{Autumn 2020}\\
				Lecturer: & \textbf{Paul Vesey \footnotesize{BEng, MIE, HDip}}\\
				Filename: & \footnotesize{\currfilename}\\
				\hline
			\end{tabular}
		\end{large}			
	}
\end{table}
\vspace{0.25cm}	
	
%     _            _                                  _     _____ 
%    / \   ___ ___(_) __ _ _ __  _ __ ___   ___ _ __ | |_  |___ / 
%   / _ \ / __/ __| |/ _` | '_ \| '_ ` _ \ / _ \ '_ \| __|   |_ \ 
%  / ___ \\__ \__ \ | (_| | | | | | | | | |  __/ | | | |_   ___) |
% /_/   \_\___/___/_|\__, |_| |_|_| |_| |_|\___|_| |_|\__| |____/ 
%                    |___/                                        

\newpage
\setcounter{page}{1}
\begin{center}
	\begin{figure}[ht]
		\centering
		\includegraphics[width = 6cm]{img/LITlogo.jpg}
		\label{fig:logoa3}
	\end{figure}
	\Large\textbf{Practical 3 - Chair and Table}\\
	\large\textbf{BA in Interior Design \& Technology}
\end{center}
In this assignment you are going to create a simple table scene, comprising table with tablecloth, chair, and set dinner place.  You may use assets from Turbosquid and other sites.\\

The scene is to be a simple contemporary table and chair (single) overlooking or beside a window.  The external scene does not need to be modelled, but the window itself is to be a prominent feature of the final image.\\

The final image is to be of photo-realistic standard, so it will be vital to keep the scene simple, and spend considerable time on lighting and materials.\\ 

In order to complete this assignment, I suggest the following workflow.  
\begin{enumerate}
	\item Model the scene in Revit
	\item Import model to 3DS and apply materials for initial render
	\item Add other 3DS objects to the scene, and complete the modelling
	\item Create and/or modify materials within another instance of 3DS.  You should perfect your materials by testing on a small object.
	\item Apply your new materials to the final model.
	\item Test camera and lighting using a variety of test renders
	\item Create your final render with at least circa 4 bounces, final gather, GI etc.
\end{enumerate}

Your final image is to be 1920 x 1024 in dimensions.  Please note that this will take considerable time to render.

Submission is to be a zip file containing the model geometry, materials, and rendered images.\\

Upload your model to Moodle upon completion.

\end{document}